\chapter{Tipps für die Durchführung einer wissenschaftlichen Arbeit}
\label{cha:Hinweise}

Die Bearbeitung einer wissenschaftlichen (Abschluss-)Arbeit gliedert sich in vier Phasen, die im Folgenden anhand von Stichworten umrissen werden.

\section{Phase 1: Einarbeitung und Literaturrecherche}
\label{sec:Einarbeitung}

\begin{itemize}
  \item dient zur Verdeutlichung der gegebenen Problemstellung
  \item führt den Bearbeiter auf den \glqq{}State-of-the-Art\grqq{} hin $\rightarrow$ Literatur
  \item zeigt Möglichkeiten zur Problemlösung auf, die dann in Phase 2 genutzt werden können
  \item dient der Ideenfindung
  \item Erstellen eines Zeitplanes
  \item regelmäßige Treffen mit dem Betreuer
\end{itemize}



\subsection*{Literaturrecherche und -management}
\label{sec:Literaturrecherche}
Zur Sammlung und Verwaltung von Literatur dienen sog. Literaturdatenbanken.
Diese Programme ermöglichen oft die direkte Verwendung als Literaturverzeichnis im späteren Text.

Das Suchen und Lesen der aktuellsten Literatur ist wichtig für eine gute Übersicht über die Problemstellung.
Über die ULB ist es möglich viele Online-Quellen%
\footnote{Für einige ist es nötig, sich im TUD-Netz zu befinden, ggf. per VPN.
Hierzu finden sich weitere Infos auf der TUD-Website.}
zu nutzen, um Texte auch als pdf-Dokumente erhalten zu können.
Die für unseren Bereich wichtigsten Seiten sind:
\begin{itemize}
  \item \href{https://scholar.google.de/}{Google Scholar}
  \item \href{http://rzblx1.uni-regensburg.de/ezeit/}{Elektronische Zeitschriftenbibliothek} (DB aller Zeitschriften und Berechtigungen)
  \item \href{http://ieeexplore.ieee.org}{IEEEXplore} (Zugang zu allen Medien der IEEE)
  \item \href{http://www.sciencedirect.com}{Sciencedirect} (Zugang zu Zeitschriften aus dem ELSEVIER Verlag)
  \item \href{http://citeseerx.ist.psu.edu}{CiteSeer} (freie DB mit Zugang zu vielen Volltexten)
  \item \href{http://www.ifac-control.org/publications/ifac-papersonline.net}{IFAC-PapersOnline} (Archiv mit Artikeln aller IFAC Konferenzen)
  \item \href{http://link.springer.com/}{Springer E-Books} (Archiv des Springer Verlags mit vielen eBooks und Zeitschriften)
\end{itemize}
auch mal auf den Webseiten der Autoren schauen und natürlich \glqq{}googlen\grqq{}.


\section{Phase 2: Bearbeitung des Problems}
\label{sec:Bearbeitung}

\begin{itemize}
  \item dient zur Lösung der Problemstellung
  \item kreativste und anstrengendste Phase der Arbeit
  \item Verwendung professioneller Hilfsmittel (Programme wie Matlab oder Mathematica etc.)
  \item Treffen wenn Bedarf besteht, keine Regelmäßigkeit mehr
  \item sowohl die Implementierung als auch die Ausarbeitung/Präsentation sollten spätestens ab dieser Phase mittels eines Versionskontrollsystems wie \bspw \textsc{Git} versioniert werden und regelmäßig zur Datensicherung auf einen externen Speicher oder einen Cloud-Speicher (\zB Hessenbox) übertragen werden, damit ältere (möglicherweise noch funktionsfähige) Stände wiederhergestellt werden können und keine Daten bei Zerstörung oder Verlust des Arbeitsrechners verloren gehen\\
    \textbf{Beachten Sie:}
    Sie als Bearbeiter der Arbeit sind für die Datensicherung selber verantwortlich.
    Dies gilt auch, wenn Sie an einem Rechner des Instituts (Rechnerpool oder Versuchstand) arbeiten!
  \item Eine gründliche Dokumentation von Teilergebnissen (Herleitungen, Rechenwege, verwendete Nomenklatur, \zB bereits in Latex) erleichtert später das Verfassen der Arbeit.
\end{itemize}


\section{Phase 3: Ergebnisse zusammenstellen}
\label{sec:Ergebnisse}

\begin{itemize}
  \item dient der Reorganisation der Arbeit
  \item sämtliche Ergebnisse werden festgehalten
  \item Strukturierung und Gliederung der Ergebnisse, so dass die nächste Phase (Schreiben) gut durchgeführt werden kann
  \item wenige, längere Treffen zur Ergebnisbesprechung mit dem Betreuer
\end{itemize}


\subsection*{Ergebnissicherung}
\label{sec:Ergebnissicherung}
\begin{itemize}
  \item alles zusammentragen was erreicht wurde $\rightarrow$ guter Überblick notwendig (s.\,o.)
  \item auch Programmcode ist ein Ergebnis $\rightarrow$ verständlich kommentieren (Englisch)
  \item auf  Wiederverwendbarkeit von Grafiken achten (Linienstärke, Farbe, Beschriftung, \ldots)
  \item nur noch kleine Änderungen durchführen (z.B. Parametereinstellungen)
\end{itemize}
\emph{Ergebnisse sollten für sich sprechen und für jeden verständlich sein (ohne Erklärung)}


\section{Phase 4: Dokumentation und Präsentation}
\label{sec:Dokumentation}

\begin{itemize}
  \item Verfassen der Arbeit und Erstellen der Präsentation
  \item letzte Verfeinerungen an den Ergebnissen (falls notwendig)
  \item schwierigste Phase der Arbeit
\end{itemize}


\subsection*{Wissenschaftliches Schreiben}
\label{sec:Wissenschaftliches Schreiben}
\begin{itemize}
  \item nicht am Layout der Arbeit aufhalten $\rightarrow$ Vorlage verwenden
  \item sehr aufwendiger Prozess von Schreiben -- Verbessern -- neu Schreiben -- \ldots
  \item Aufwand darf nicht unterschätzt werden (Richtwert: 1-2 Seiten/Tag)
  \item einfach erst einmal aufschreiben --  korrigiert wird dann später
  \item Struktur einer wissenschaftlichen Arbeit ist vorgegeben
    \begin{itemize}
      \item \textbf{Titelseite} mit Art der Arbeit, Titel, Namen des Autors sowie Abgabedatum.
      \item \textbf{Aufgabenstellung} wird vom Betreuer der Arbeit zur Verfügung
        gestellt.
      \item \textbf{Erklärung} zur Selbständigkeit.
        Der Text ist vorgegeben und wird bei Verwendung der Vorlage automatisch erzeugt.
      \item \textbf{Kurzfassung} der Arbeit.
        Der Umfang soll so bemessen sein, dass die englische Version (\textbf{Abstract}) auf die gleiche Seite passt.
      \item \textbf{Inhaltsverzeichnis} wird in \LaTeX{} durch \verb|\tableofcontents| automatisch erzeugt.
      \item \textbf{Symbole und Abkürzungen}.
        Dieses Verzeichnis erstellt man am Besten von Hand.
        Die Einteilung in lateinische und griechische Symbole und Formelzeichen kann nach Bedarf geändert werden (zum Beispiel nach Kapiteln oder Konzepten) oder ganz weggelassen werden.
      \item \textbf{Hauptteil der Arbeit}, in einzelne Kapitel und Abschnitte unterteilen.
      \item \textbf{Anhang}.
        Hier können Abschnitte stehen, die beim Lesen der Arbeit stören würden, \zB\ Programmcode, technische Daten oder lange mathematische Beweise.
      \item \textbf{Literaturverzeichnis} sollte automatisch generiert (in \LaTeX{} \zB mit biblatex/biber). Formatierung und Inhalt prüfen und ggf. korrigieren!
      \item \label{itm:Ordner} \textbf{Zusätzliches Material} wie \zB\ der vollständige Programmcode eines Software-Projekts gehört nicht in die Arbeit, sondern kann in einem separaten Ordner abgelegt werden.
    \end{itemize}
  \item Wir bieten Ihnen an, am Ende einmal Ihre Arbeit korrektur zu lesen.
    Darüberhinaus können Sie sich mit konkreten Fragen zu einzelnen Aspekten natürlich auch immer an Ihren Betreuer wenden.
  \item Orientieren Sie sich beim Schreiben an "`POEM"':\\
    $P_\text{rägnanz}\; O_\text{rdnung}\; E_\text{infachheit}\; M_\text{otivation}$ (einfache, kurze, strukturierte und anregende Sätze)
\end{itemize}


\paragraph{Inhalt und Umfang der Arbeit}
Ihre Arbeit soll am Ende \emph{vollständig} sein, und dies hat zunächst nichts mit der Länge in Seiten zu tun, sondern mit inhaltlichen Fragen:
\begin{itemize}
  \item Wird die Problemstellung und Voraussetzungen der Arbeit klar?
  \item Kann man Ihrem Lösungsansatz folgen?
  \item Wird das Ergebnis ausreichend diskutiert und eingeordnet, \zB durch Vergleiche mit alternativen Ansätzen oder bestehenden Methoden?
  \item Wird bei negativen Ergebnissen eine überzeugende Analyse durchgeführt, warum ein gewählter Ansatz nicht funktioniert hat?
\end{itemize}
Je nach Art der Arbeit umfasst eine vollständige Angabe noch weitere Aspekte, die nicht unbedingt Teil der eigentlichen Ausarbeitung sind:
\begin{itemize}
  \item Sind (bei konstruktiven Arbeiten) die Konstruktionsunterlagen vollständig?
  \item Ist Programmcode, der weiterverwendet werden soll, entsprechend ausführlich kommentiert?
  \item Gibt es zu dem entwickelten Versuchsaufbau eine verständliche Anleitung?
\end{itemize}

Es soll an dieser Stelle bewusst kein erwarteter Umfang der Arbeit in Seiten angegeben werden, da dieser je nach der Art der Arbeit (Bachelorarbeit, Masterarbeit, \ldots) und Ausprägung (theoretisch, simulativ, experimentell, konstruktiv) der Arbeit zu stark schwankt.
Bei Fragen und Unsicherheiten zum Umfang der Arbeit wird Ihnen daher Ihr Betreuer weiterhelfen.





\subsection*{Verteidigung und Präsentation}
\label{sec:Verteidigung}
\begin{itemize}
  \item Dauer der Präsentation (ohne Fragerunde): 20 Minuten (bei Bachelor- und Masterarbeiten)
    \begin{itemize}
      \item nach 25 Minuten wird abgebrochen
      \item Grundsätzlich lieber etwas zu lang (aber auf jeden Fall unter 25 Minuten) als zu kurz.
      \item Bei Proseminaren (10 Minuten) und Projektseminaren (je nach Anzahl der Teilnehmer) abweichende Länge
    \end{itemize}
  \item Inhalt der Arbeit auf ca. 10 wesentliche Punkte reduzieren
  \item je Punkt eine Folie, maximal zwei
  \item je Folie 1-2 min Gesprächszeit
  \item Struktur der Präsentation:
  \begin{itemize}
    \item Motivation / Einleitung (Aufgabenstellung)
    \item Grundlagen
    \item Lösungsweg
    \item Ergebnisse
    \item Zusammenfassung (und Ausblick)
  \end{itemize}
  \item Beachten Sie: Was sich durch ein Bild (anstelle von Text) darstellen lässt, lässt sich durch den Zuhörer schneller erfassen und wertet die Präsentation optisch auf.
  \item insbesondere komplexe Sachverhalte durch Abbildungen verdeutlichen
  \item nur Stichpunkte schreiben, keine ganzen Sätze
  \item klare, einfach nachvollziehbare Notation (\zB bei Variablen) verwenden
  \item mehrmaliges Üben des Vortrags (\zB vor dem Spiegel, vor der Familie, \ldots)
  \item Probevortrag vor Betreuer halten
\end{itemize}
\emph{nichts weglassen was auf einer Folie steht, gerne zusätzliche Dinge erwähnen}


\section{Allgemeine Punkte}

\subsection*{Organisatorisches}
\label{sec:Organisatorisches}

\begin{itemize}
  \item jeder Studierende sollte an mindestens einem Regelungstechnischen Seminar teilnehmen
  \item Abgabetermin ist ein fixer Termin $\rightarrow$ Prüfungsleistung (Durchfallen möglich)
  \item jeder Studierende erhält auf Wunsch einen eigenen Account für unsere Rechner (bei Betreuer melden)
  \item zur Erstellung der Ausarbeitung und der Präsentation wird die Verwendung des Textsatzprogramms \LaTeX{} empfohlen, das Template für die Ausarbeitung erhalten Sie von Ihrem Betreuer
\end{itemize}




\subsection*{Generelles zu Treffen mit dem Betreuer}
\label{sec:TreffenBetreuer}

Die persönlichen Treffen mit dem Betreuer sollten gut vorbereitet werden (\zB kurze Präsentation, Fragenkatalog).
Die Ergebnisse der Treffen und die Anmerkungen des Betreuers sollten Sie in Form von Notizen festhalten.


Hier ist zu bemerken, dass ein häufigeres Treffen mit dem Betreuer alleine keine mangelnde Selbständigkeit impliziert.
Solange diese Treffen derart ablaufen, dass Sie Ihren Arbeitsfortschritt seit dem letzten Treffen auf \zB ein paar kurzen Folien vorstellen (die nicht übermäßig "`schön"' gestaltet sein sollen), Probleme benennen, die dabei möglicherweise aufgetreten sind und dann noch erläutern, wie Sie mit diesen umgehen wollen, wird es auch bei einem wöchentlichen Treffen hieraus keine Abzüge in der Selbständigkeit geben.
Wenn Sie dagegen der Auf"|fassung sind, Ihre Selbständigkeit dadurch unter Beweis zu stellen, dass Sie monatelang nichts von sich hören lassen, ist dies hinsichtlich der Kommunikation mit dem Betreuer nicht ideal, und vor allem tragen Sie dann das volle Risiko, wenn Sie die Aufgabe falsch verstanden haben und damit in die falsche Richtung arbeiten!




\subsection*{Allgemeines zu Hilfsmitteln bei der Erstellung und Bearbeitung}
\label{sec:Allgemeines}

\begin{itemize}
  \item es gibt viele Bücher die bei der Erstellung von wissenschaftlichen Arbeiten helfen
  \item viele verwendete Programme bieten ausführliche Hilfen an (\zB \Matlab, Mathematica)
  \item bei Problemen immer zuerst in der Hilfe schauen, dann "`googlen"', dann Betreuer fragen
\end{itemize}
